% ---------------------------------------------------------------------------
\section*{Key findings}
\addcontentsline{toc}{section}{Key findings}

The first seven findings are, to our knowledge, \textbf{new}; the last three strengthen earlier ideas with new
measurements.

\begin{enumerate}[label=\textbf{\arabic*.}]
\item \textbf{The text is a chain of glyphs in which the space is a weak boundary.} The rules linking the last glyph of
  a word to the first glyph of the next mirror the rules between two glyphs inside a word more than in any of 71
  comparison languages. Spaces can be predicted from the two surrounding glyphs (accuracy 0.86, above 97\% of
  languages), and cutting the text at other points allowed by the rules mostly yields words attested elsewhere in the
  manuscript (50\%, against 6.5\% in languages).
\item \textbf{Adjacent words are linked by junction rules.} \eva{qo-} follows words ending in \eva{-y}, \eva{-o},
  \eva{-d}; \eva{o-} follows \eva{-n}, \eva{-r}, \eva{-s}, \eva{-m}. A final \eva{-l} or \eva{-r} depends on how the
  next word begins. The scribe \emph{knows the next word} while finishing the current one. The link is at least four
  times stronger than in any medieval scribe we tested.
\item \textbf{The line is a closed unit.} The link between words stops at the line break and also where a drawing
  interrupts the line; line starts and ends have their own forms (at line end \eva{-m} replaces \eva{-r}; at line start
  \eva{s-} and \eva{y-} appear).
\item \textbf{The scribe copies from the line immediately above.} Words reproduce, more than chance, words or pieces of
  line from the line just written, not from lines further up; the copy is adapted to the junction rules with its new
  neighbour.
\item \textbf{Each line avoids starting like the line above.} If the line above begins with \eva{qo-}, the line below
  begins with \eva{o-}. None of 79 medieval scribes does this: not a single case in 72 tests on the six Nordic scribes,
  and only three of 73 German manuscripts avoid one single initial, as expected by chance.
\item \textbf{Choices between variant forms follow a ``state'' lasting a few words.} Between the forms \eva{k/t},
  \eva{sh/ch}, \eva{-ey/-dy}, the scribe tends to repeat the same choice for two or three words, \emph{whatever word is
  being written}. It is not copying of similar words, not position in the line, it is separate for each choice and
  identical in all the hands of the manuscript: a rule of the writing system. We did not find it in the six Nordic
  scribes nor in the Copiale copyist (the test on the 73 German scribes is in progress).
\item \textbf{Marked variants drift along the line} (\eva{qo}, \eva{k}, \eva{sh}, \eva{-ey} become rarer from left to
  right), three to fifteen times more than in the six Nordic scribes (German scribes in progress).
\item \textbf{The Voynich repeats the previous word immediately} (1.0\% of adjacent pairs) or a nearly identical one
  (3.9\%), more than every language and every one of the 85 handwritten texts we tested (usually 10 to 100 times more),
  while having a vocabulary of normal size.
\item \textbf{The book has a writing structure}: each page has a strong identity, the bifolium is a working unit (one
  Currier ``language'' and almost always one hand), the paragraph opens with a large sign (\eva{p}) and has first and
  last lines with their own register; \eva{-ey} increases down the page.
\item \textbf{No proposed reading survives the controls}: neither ours (substitutions in dozens of languages, solvers,
  verbose ciphers, nomenclators, content-guided searches) nor the published ones we could test (Bax, Vatne, Cheshire,
  Schechter, Gatta), nor the historical cipher mechanisms we simulated (homophones, Alberti's index, Bacon's biliteral
  cipher, Trithemius' word lists, autokey).
\end{enumerate}

% ---------------------------------------------------------------------------
\section{Introduction}

\subsection{The manuscript}

The Voynich manuscript is a codex of about 240 pages, with drawings of plants, astronomical and zodiacal diagrams,
female figures in pools and pipes, pharmacy jars, and a long final section of text only, divided into short paragraphs
marked by stars (usually called ``recipes''). The script is neat and confident: it does not look like the work of
someone inventing signs while writing. In the 1970s Prescott Currier noticed that the pages fall into two varieties of
writing, now called \textbf{language A} and \textbf{language B}, and that more than one hand wrote the manuscript;
Lisa Fagin Davis later distinguished five scribes.

The text is studied through \textbf{transliterations} that render each sign with a conventional Latin letter. The most
used is the EVA alphabet: a word such as \eva{qokeedy} or \eva{daiin} says nothing about sound; it is only a label for a
sequence of signs. Some signs are written in EVA with more than one letter (\eva{ch}, \eva{sh}, the tall ``gallows''
\eva{k}, \eva{t}, \eva{p}, \eva{f} and their compounds); in all our measurements we count them as a single sign.

\subsection{The problem}

Over more than a century dozens of readings have been proposed: abbreviated Latin, Romance languages, Hebrew, Old
Turkic, Asian languages, ciphers of various kinds. None has convinced the scholarly community, for a simple reason: it
is easy to find in a long text a few words that ``look like'' something, and hard to show that a reading does better
than chance. On the other side, those who argue that the text has no meaning must explain why it is so regular.

\subsection{What this work offers}

Not a new reading, but a set of \textbf{measurements with controls}:
\begin{itemize}
\item each measurement was \textbf{written down before} looking at the result on the Voynich (preregistration);
\item each method was \textbf{first tested on synthetic texts} in which the effect was present or absent, to check that
  it recovers it;
\item each property of the Voynich was \textbf{compared} with the same measurement on languages, scribes, ciphers and
  generators;
\item when an intermediate measurement proved wrong, \textbf{we wrote it down} and corrected it (the internal checks are
  listed in Appendix~B).
\end{itemize}
The result is a description of the Voynich as a \textbf{writing system}: what the writer does, word after word and line
after line. This description is useful in two ways: it says which explanations are excluded, and it provides a list of
properties that any future explanation will have to reproduce.

% ---------------------------------------------------------------------------
\section{Data}

\subsection{The Voynich}

Three independent transliterations in René Zandbergen's IVTFF format:

\begin{center}\small
\begin{tabular}{lll}
\toprule
transliteration & alphabet & role \\
\midrule
Zandbergen--Landini (ZL) & EVA & reference \\
Takahashi (IT) & basic EVA & replication \\
Glen Claston (GC) & v101 (a different segmentation of the signs) & replication \\
\bottomrule
\end{tabular}
\end{center}

\noindent We use the paragraph text, about 35{,}000 legible words. Labels of drawings, circular texts and radii enter
only some tests, which are indicated. Words with illegible signs (0.6\%) are discarded. Every main result was replicated
on at least two transliterations, most on all three.

\subsection{Comparison texts}

\begin{itemize}
\item \textbf{Natural languages:} the Bible in about 100 languages; 71 texts from the Gaskell and Bowern (2022)
  collection, as whole words; Latin technical texts (Pliny, Varro, Isidore, Vitruvius) and an Italian one (Alberti's
  \emph{Della Pittura}), to control for genre.
\item \textbf{Constructed languages:} Esperanto, Volapük, Neo-Quenya, Interlingua, Lojban, Klingon, LOLCat, Toki Pona.
\item \textbf{Handwritten gibberish:} 38 texts produced by volunteers asked to write an invented text by hand (Gaskell and
  Bowern 2022).
\item \textbf{Published generators and ciphers that imitate the Voynich:} the Naibbe cipher (Greshko 2025), which
  encrypts a real Latin or Italian text with a procedure executable by hand; the ``self-citation'' generator of Timm and
  Schinner (2020), which produces text without a message by copying and modifying words already written; two amateur
  generators (U2, U3).
\item \textbf{Real medieval scribes, at the level of the facsimile} (that is, with letter forms as in the manuscript,
  abbreviations, and the original lines and pages):
  \begin{itemize}
  \item six manuscripts from the \textbf{Menota} archive (Medieval Nordic Text Archive): Icelandic, Norwegian and one
    Swedish, from about 1200 to 1550;
  \item 73 German manuscripts from the \textbf{Referenzkorpus Frühneuhochdeutsch} (ReF), from 1350 to 1500, that is, from
    the area and century of the Voynich, with lines of 5--8 words like its own;
  \item the Anglo-Saxon scribe of the \emph{Hatton Gospels} (from an edition, without original lines).
  \end{itemize}
\item \textbf{A real historical cipher:} the \textbf{Copiale}, an eighteenth-century German enciphered manuscript (105
  pages), with the transcription and decipherment of Knight, Megyesi and Schaefer (2011). It uses a homophonic cipher:
  the same letter can be written with different symbols, at the copyist's choice.
\item \textbf{Handwritten medieval technical texts:} from the ReF, a \emph{Book of Nature}, a surgical recipe book, a
  book of lunar astrology, a \emph{Naturlehre} and a book of magic and alchemy.
\end{itemize}
All data were taken at fixed versions, with the licences listed in Appendix~D. Protected texts are not redistributed.

% ---------------------------------------------------------------------------
\section{Methods}

\subsection{The rules we set ourselves}

\begin{enumerate}
\item \textbf{Preregistration.} For each experiment, before running it on the Voynich, we wrote in a dated file the
  question, the measure, the data and the \textbf{criterion} by which we would read the outcome (``if the value exceeds
  $X$ we will say $Y$'').
\item \textbf{Testing the code on synthetic texts.} Every measure was first tested on texts built for the purpose, with
  the effect present and absent. A measure that did not distinguish the two cases was not used: in three cases the
  experiment was not run and was declared as a gap (Section~11).
\item \textbf{Positive and negative controls.} For decipherment tests: a known text enciphered in the same way (which the
  method must solve) and a text in another language (which shows what is obtained by chance).
\item \textbf{Honest intervals.} Uncertainty intervals are computed by resampling whole pages, because neighbouring words
  are not independent.
\item \textbf{A failed step is repeated once}, then the gap is declared.
\item \textbf{Corrections are written down.} The laboratory notebook was written only by adding, never by deleting.
\end{enumerate}

\subsection{The main measures, in plain words}

\begin{itemize}
\item \textbf{Junction:} how much information the last glyph of a word gives about the first glyph of the next word,
  beyond what glyph frequencies alone would give (in bits).
\item \textbf{Junction rule:} how much a choice at the start (or end) of a word depends on the glyph of the neighbouring
  word, for the same word.
\item \textbf{Agreement of choices (the ``state''):} for a choice between two variants (for example \eva{k} or \eva{t}),
  how much two nearby words of the same line make the same choice more than expected, where the expectation accounts for
  the preferences of each word, of the page and of the position in the line. It is expressed as a fraction of the
  maximum possible agreement (0 = chance, positive values = more agreement than chance), like Cohen's kappa.
\item \textbf{Copying from the line above:} how many words have an identical or nearly identical word in the line above,
  beyond chance.
\item \textbf{Left margin:} how often a line begins like the line above, compared with shuffled lines.
\end{itemize}

\subsection{A decisive methodological correction}

Midway through the work we found that the measure of agreement between choices had a \textbf{bias} that varied from text
to text: in some texts it inflated agreement by up to $+0.37$, in others it lowered it. The cause is that the expectation
is estimated from the same data. The correction consists in subtracting the value obtained by shuffling the choices among
occurrences of the same word, which has the same bias but no real agreement. All agreement measures reported here are
corrected in this way. The correction changed three provisional indications of the work: the Anglo-Saxon scribe does not
have the ``window'' a first measure had attributed to him; Volapük does not resemble the Voynich; and the Voynich does
not lie \emph{outside} the range of languages for the shape of its agreement, but at its top.
